\chapter*{Abstract}

\vspace{1em}

{\centering
  \fontsize{22pt}{22pt}\selectfont An Interpretation of the St. Petersburg Paradox \\ using a Trial-Probability of Profit Model
  \par}

\vspace{2em}

\begin{center}
  \begin{minipage}{0.8\textwidth}
    {\raggedleft
      \fontsize{14pt}{14pt}\selectfont
      Gunwoo Kim\\
      Department of Computer Science and Engineering\\
      The Graduate School\\
      Seoul National University
      \par}
  \end{minipage}
\end{center}

\vspace{1em}

\begin{spacing}{1.2}
  \begin{center}
    \begin{minipage}{0.8\textwidth}
      This study proposes a trial-probability of profit model to address the St. Petersburg Paradox and analyze human decision-making processes. Unlike traditional approaches relying on utility functions or external constraints, this research focuses on the probability of profit over repeated trials. The findings reveal that the paradox stems not from infinite expected value but from the slow convergence in probability of profit under sparse reward structures. Experimental results demonstrate that higher participation costs further decelerate convergence in probability of profit, explaining why individuals deem high cost games irrational. The proposed model offers a novel framework for interpreting decision-making in scenarios involving low probability, high reward outcomes, providing insights into the St. Petersburg Paradox.
    \end{minipage}
  \end{center}
\end{spacing}

\vspace{1em}

keywords : St. Petersburg paradox, decision-making, expected utility
